\documentclass[10pt,a4paper]{amsart}
\usepackage[margin=19mm]{geometry}
\usepackage{amsmath,amssymb,mathtools}
\usepackage[scr=rsfs,cal=euler]{mathalfa}
\usepackage{xfrac}
\usepackage[unicode,hidelinks,bookmarks=false]{hyperref}
\newcommand{\ie}{i.e.\ }
\newcommand{\cf}{cf.\ }
\newcommand{\resp}{respectively\ }
\newcommand{\Subsection}{\subsection}
\providecommand{\nobreakdash}{\nobreak\hbox{-}}
\setlength{\emergencystretch}{2em}
\multlinegap0pt
\usepackage{amssymb}
\usepackage{dsfont}
\usepackage{xcolor}
\newtheorem{lemma}{Lemma}
\title{Global universality of the expected number \\of zeros of non-analytic random signals}
\author{J\"urgen Angst, Thibault Pautrel, Guillaume Poly}
\date{Source: arXiv:2609.01007v1 (CC BY 4.0)}
\begin{document}
\maketitle

\noindent\emph{Source and licence.} Global universality of the expected number \\of zeros of non-analytic random signals. Version arXiv:2609.01007v1 (CC BY 4.0), distributed by arXiv under the Creative Commons Attribution 4.0 International licence, http://creativecommons.org/licenses/by/4.0/, as recorded in the arXiv licence metadata for this exact version and shown on the abstract page https://arxiv.org/abs/2609.01007v1. Full licence notice: \texttt{licenses/arxiv-2609.01007-CC-BY-4.0.txt}.\par\smallskip

\noindent\textit{Editorial scope.} Counted unit: the article's lemma lem.decorel (source0.tex lines 824--828). The article's complete original proof of it is delivered with it (source0.tex lines 829--887, one \texttt{proof} environment of 59 lines). No mathematics is written, repaired or re-derived: the statement and the proof are the article's own lines, in the article's own notation, and no retained line is substituted or re-flowed.  No further source-local context is needed: every cross-reference of every retained line resolves inside the two delivered spans.  Nothing was omitted from any retained span, so the package carries no omission record.  No retained line contains a citation, so the wrapper carries no bibliography and no bibliography entry.  No retained line contains a figure environment, an image-inclusion command or drawing code, so no image is delivered and the package declares no asset dependency.  Editorial formatting only, in addition: an AMS \texttt{amsart} preamble that restates the article's own declarations and macros used by the retained lines, namely the environments \texttt{lemma} on separate counters, together with the standard packages those lines need.  The retained environments print in this document as Lemma 1 in order of appearance, whereas the article prints the same items with the numbering of its own full document.  The complete native source of the article as distributed in the arXiv e-print for this exact version is bundled unchanged as \texttt{sources/209073/source0.tex}.
\begin{lemma}{Suppose that the base function $f$ satisfies the condition $\mathbf{(H1)}$}, then there exists a constant $C=C(||f||_{H^1})$ such that
\[
\mathbb E_{X,Y} \left[ \left|\mathbb E[f_n(X)f_n(Y)] \right|\right]\leq \frac{C}{\sqrt{n}}.
\]\label{lem.decorel}
\end{lemma}

\begin{proof}[Proof of Lemma \ref{lem.decorel}]
We have
\[
\mathbb E[f_n(X)f_n(Y)]=\frac{1}{n}\sum_{k,\ell=1}^{n}\mathbb E[a_ka_{\ell}]f(kX)f(\ell Y)=\frac{1}{n}\sum_{k=1}^{n}f(kX)f(kY).
\]
{Since the base function $f$ satisfies the condition $\mathbf{(H1)}$, it is the limit of its Fourier sums $S_N f(x):=\sum_{|p| \leq N} \widehat{f}(p) e^{i p x}$, and since it is centered, we have $ \widehat{f}(0) =0$}. For any $N$, we can then write 
\[
\begin{array}{ll}
\mathbb E[f_n(X)f_n(Y)]  =\displaystyle{\frac{1}{n}\sum_{k=1}^{n} \left( f(kX)-S_N f(kX)\right) f(kY)}\\
\\
  + \displaystyle{\frac{1}{n}\sum_{k=1}^{n} S_N f(kX) \left(f(kY)-S_N f(kY)\right)
  +\frac{1}{n}\sum_{k=1}^{n} S_N f(kX) S_N f(kY).}
\end{array}
\]
Now taking the absolute value and the expectation under $\mathbb E_{X,Y}$, we get 
\[
\begin{array}{l}
\mathbb E_{X,Y} \left[ \left|\mathbb E[f_n(X)f_n(Y)] \right|\right] \leq \displaystyle{\frac{1}{n}\sum_{k=1}^{n} \mathbb E_X\left[ \left| f(kX)-S_N f(kX)\right|\right] \mathbb E_Y \left[\left|f(kY)\right|\right]}\\
\\
  + \displaystyle{\frac{1}{n}\sum_{k=1}^{n} \mathbb E_X\left[|S_N f(kX)| \right] \mathbb E_Y\left[ |f(kY)-S_N f(kY)|\right]} 
+ \displaystyle{\mathbb E_{X,Y}\left[ \left| \frac{1}{n}\sum_{k=1}^{n} S_N f(kX) S_N f(kY)\right|\right].}
\end{array}
\]
Applying Cauchy--Schwarz inequality, we get for the first term on the right hand side
\[
\frac{1}{n}\sum_{k=1}^{n} \mathbb E_X\left[ \left| f(kX)-S_N f(kX)\right|\right] \mathbb E_Y \left[\left|f(kY)\right|\right] \leq   ||f-S_Nf ||_2 \times ||f||_2.
\]
The second term can be treated similarly, namely 
\[
\frac{1}{n}\sum_{k=1}^{n} \mathbb E_X\left[|S_N f(kX)| \right] \mathbb E_Y\left[ |f(kY)-S_N f(kY)|\right]\leq \, ||f||_2 \times  ||f-S_Nf ||_2.
\]
For the third term, {recalling that $ \widehat{f}(0)=0$}, observe that 
\[
S_Nf(kX) S_Nf(k Y)
= \sum_{1\leq |p|\leq N,1\leq |q|\leq N}\widehat{f}(p)\widehat{f}(-q)e^{ik(pX-qY)} 
\]
and
$$\frac{1}{n}\sum_{k=1}^{n}e^{ik(pX-qY)}=\exp\left(i\frac{n+1}{2}(pX-qY)\right)\frac{1}{n}\frac{\sin\left(\frac{n}{2}(pX-qY)\right)}{\sin\left(\frac{pX-qY}{2}\right)}.$$
Therefore, applying once again Cauchy--Schwarz inequality (twice indeed), we deduce that if $K_n$ denotes the standard Fej\'er such that $||K_n||_1=1$
\[
\begin{array}{l}
\displaystyle{\frac{1}{n}\sum_{k=1}^{n} \mathbb E_{X,Y}\left[|S_N f(kX) S_N f(kY)|\right]}  \displaystyle{\leq \sum_{\substack{1\leq |p|\leq N, \\ 1\leq |q|\leq N}}|\widehat{f}(p)\widehat{f}(-q)|\mathbb E_{X,Y}\left[\frac{1}{n}K_n(pX-qY)\right]^{1/2}}\\
\\
 =\displaystyle{ \frac{1}{\sqrt{n}} \left( \sum_{1\leq |p|\leq N}|\widehat{f}(p)|\right)^2 \leq \frac{1}{\sqrt{n}} \left( \sum_{1\leq |p|\leq N}|\widehat{f'}(p)|^2\right)\left( \sum_{1\leq |p|\leq N}\frac{1}{p^2}\right)\leq \frac{2 ||f'||_2^2}{\sqrt{n}}}.
\end{array}
\]
We thus get 
\[
\mathbb E_{X,Y} \left[ \left|\mathbb E[f_n(X)f_n(Y)] \right|\right] \leq 2\, ||f||_2 \times ||f-S_Nf ||_2 + \frac{2 ||f'||_2^2}{\sqrt{n}}.
\]
Now by Parseval identity, we have 
\[
||f-S_Nf ||_2^2 = \sum_{|p|>N}|\widehat{f}(p)|^2 =  \sum_{|p|>N}|\widehat{f'}(p)|^2 \times \frac{1}{p^2} \leq \frac{||f'||_2^2}{N^2}.
\]
As a conclusion, choosing the threshold $N=\lfloor\sqrt{n}\rfloor$, we get indeed that there exists a constant $C=C(||f||_{H^1})$ such that 
\[
\mathbb E_{X,Y} \left[ \left|\mathbb E[f_n(X)f_n(Y)] \right|\right]\leq \frac{C}{\sqrt{n}}.
\]
\end{proof}

\end{document}
